\section{Advanced Unpacking and Functions}
\label{sec:functions_unpacking}

\subsection{Tuple Packing and Unpacking}
Python allows for the grouping and distribution of objects through packing and unpacking mechanisms.

\begin{itemize}
    \item \textbf{Packing:} Multiple individual objects are grouped into a single tuple.
    \item \textbf{Unpacking:} Items within an iterable are assigned to individual variables. This requires the number of variables to match the number of elements unless extended unpacking is used.
\end{itemize}

\begin{lstlisting}[language=Python]
# Packing
employee = 'Peter', 'Muller' # Results in ('Peter', 'Muller')

# Unpacking
first, last = employee

# Multiple Assignment (combined packing/unpacking)
x, y, z = 5, 10, 15
\end{lstlisting}

\subsubsection{Extended Iterable Unpacking}
Unpacking works with any iterable (strings, lists, sets) and supports specific operators to handle variable lengths or unwanted values.

\begin{itemize}
    \item \textbf{Rest Operator (*):} Captures remaining elements into a list.
    \item \textbf{Placeholder (\_):} Used by convention to ignore specific values during unpacking.
\end{itemize}

\begin{lstlisting}[language=Python]
# Unpacking with rest
first, last, *address, birth_date = ("Peter", "Muller", "Street", 10, "14.04.1977")
# address becomes ["Street", 10]

# Ignoring values
first, last, _, birth_date = employee_data

# Unpacking in nested structures
x = [1, 2, (3, [4, 5])]
a, b, (c, (d, e)) = x
\end{lstlisting}

\subsubsection{Iterable Expansion}
The \texttt{*} and \texttt{**} operators allow unpacking iterables within other collection literals.

\begin{lstlisting}[language=Python]
list1 = [1, 2]
list2 = [*list1, 3, 4] # [1, 2, 3, 4]

dict1 = {"A": 1, "B": 2}
dict2 = {"C": 3, **dict1} # {'C': 3, 'A': 1, 'B': 2}
\end{lstlisting}

\subsection{Functions}
Functions are defined using the \texttt{def} keyword. They establish a local scope and can return objects to the caller.

\begin{itemize}
    \item \textbf{Parameters:} Placeholders defined in the function signature.
    \item \textbf{Arguments:} Actual values passed during the call.
    \item \textbf{Return Value:} Functions return \texttt{None} by default. Multiple values can be returned as a tuple.
\end{itemize}

\subsubsection{Argument Types and Constraints}
\begin{itemize}
    \item \textbf{Positional vs. Keyword:} Positional arguments must precede keyword arguments in a call.
    \item \textbf{Default Parameters:} Parameters with default values must follow those without.
    \item \textbf{Variadic Arguments (*args, **kwargs):} \texttt{*args} collects extra positional arguments into a 
    \crd{\textbf{tuple}}; \texttt{**kwargs} collects extra keyword arguments into a dictionary.
\end{itemize}

\subsubsection{Argument Enforcement}
Developers can force specific call styles using \texttt{/} and \texttt{*}.

\begin{lstlisting}[language=Python]
def constrained_func(pos_only, /, pos_or_kw, *, kw_only):
    pass

# / : All arguments before must be positional-only.
# * : All arguments after must be keyword-only.
\end{lstlisting}

\subsection{Docstrings}
Documentation strings should follow the function definition to describe its behavior. They can be retrieved via \texttt{help(function\_name)} or the \texttt{\_\_doc\_\_} attribute.

\begin{lstlisting}[language=Python]
def add(a, b):
    """Return the sum of a and b."""
    return a + b
\end{lstlisting}

\subsection{Lambda Functions}
Lambdas are anonymous, single-expression functions. They are primarily used as "key" functions for sorting or filtering.

\textbf{Syntax:} \texttt{lambda parameters: expression}

\begin{lstlisting}[language=Python]
# Sorting a list of strings by the numeric part
articles = ['art1', 'art20', 'art3']
sorted_list = sorted(articles, key=lambda x: int(x[3:]))
\end{lstlisting}